% ============================================================================
% BILDFILTER-PROJEKT (1f) - Main Document
% Compile with: xelatex skript.tex (run twice for correct overlays)
% ============================================================================
\documentclass[11pt, a4paper, twoside, openright]{book}

% --- Shared preamble modules ---
\input{../preamble/packages}
\input{../preamble/colors}
\input{../preamble/commands}
\input{../preamble/environments}
\input{../preamble/layout}

% --- Python code listings (matches 1f_FS26_algorithmen style) ---
\usepackage{listings}
\lstdefinestyle{python}{
  language=Python,
  basicstyle=\ttfamily\small,
  keywordstyle=\bfseries\color{maincolor},
  commentstyle=\itshape\color{maincolor!50!black},
  stringstyle=\color{headerbar},
  numbers=none,
  backgroundcolor=\color{boxbody},
  frame=single,
  rulecolor=\color{maincolor!30},
  breaklines=true,
  breakatwhitespace=true,
  tabsize=4,
  showstringspaces=false,
  aboveskip=8pt,
  belowskip=8pt,
  xleftmargin=4pt,
  framexleftmargin=4pt,
}
\lstset{style=python}

% --- Pseudocode style: visibly different from Python so students translate
%     rather than copy. Sans-serif, cinnabar accent, no Python keywords. ---
\definecolor{pseudocinnabar}{RGB}{177, 56, 46}
\lstdefinestyle{pseudo}{
  language=,
  basicstyle=\sffamily\small,
  keywordstyle=\bfseries\color{pseudocinnabar},
  morekeywords={Funktion, fuer, Fuer, FUER, jede, jeden, jedes, des, der,
                von, bis, in, an, gib, zurueck, wenn, dann, sonst, und, oder,
                nicht, ist, gleich, kleiner, groesser, leere, Liste, Setze,
                erstelle, lies, schreibe, fuege, hinzu, anhaengen, hange},
  numbers=none,
  backgroundcolor=\color{pseudocinnabar!4},
  frame=single,
  rulecolor=\color{pseudocinnabar!40},
  breaklines=true,
  breakatwhitespace=true,
  tabsize=2,
  showstringspaces=false,
  aboveskip=8pt,
  belowskip=8pt,
  xleftmargin=4pt,
  framexleftmargin=4pt,
}

% --- Script-specific bibliography ---
\addbibresource{references.bib}

% ============================================================================
% DOCUMENT
% ============================================================================
\begin{document}

\makeatletter
\def\cleardoublepage{\clearpage\if@twoside\ifodd\c@page\else
  \thispagestyle{empty}\hbox{}\newpage
  \if@twocolumn\hbox{}\newpage\fi\fi\fi}
\makeatother

% ============================================================================
% TITLE PAGE
% ============================================================================
\begin{titlepage}
\thispagestyle{empty}

\vspace*{30mm}

% --- Title ---
\begin{center}
  {\fontsize{28}{34}\selectfont\bfseries\color{headerbar}%
   Bildfilter in Python}\\[8pt]
  {\fontsize{18}{22}\selectfont\color{headerbar}%
   Projekt mit Stempel-Wanderung}
\end{center}

\vspace{10mm}

% --- Horizontal rule ---
\noindent\textcolor{headerbar}{\rule{\textwidth}{0.8pt}}

\vspace{12mm}

% --- Authors ---
\noindent{\large\bfseries Lars Rafeldt}

\vspace{6pt}

\noindent{\itshape
MNG, Mathematisch-Naturwissenschaftliches Gymnasium R\"amibuhl\\
Fachschaft Informatik}

\vspace{6pt}

\noindent{\itshape E-mail:}
{\color{maincolor}\href{mailto:lars.rafeldt@mng.ch}{lars.rafeldt@mng.ch}}

\vspace{6pt}

\noindent{\itshape Klasse 1f, Fr\"uhlingssemester 2026}

\vfill

% --- Bottom: logo spanning full page width, 5mm from page bottom ---
\begin{tikzpicture}[remember picture, overlay]
  \node[anchor=south, inner sep=0pt]
    at ([yshift=5mm]current page.south)
    {\includegraphics[width=\paperwidth]{../img/info_logo2}};
\end{tikzpicture}

\end{titlepage}

\thispagestyle{empty}
\null\newpage

\setcounter{page}{1}

% ============================================================================
% CHAPTER 1: Projektbeschreibung
% ============================================================================
\renewcommand{\chaptertag}{Bildfilter}
\chapter{Projektbeschreibung}

\vspace{-5pt}
{\color{maincolor}\rule{\textwidth}{0.4pt}}
\vspace{4pt}

{\color{maincolor}\textbf{Lernziele:}}
\begin{itemize}[leftmargin=*]
  \item Du verstehst, wie ein Computer ein Bild speichert.
  \item Du kannst eigene Bildfilter als Python-Funktionen schreiben und an
    echten Fotos ausprobieren.
  \item Du gehst Schritt f\"ur Schritt von einfachen Pixel-\"Anderungen
    bis zu Filtern, die mehrere Nachbarpixel zusammenrechnen.
\end{itemize}
\vspace{-2pt}
{\color{maincolor}\rule{\textwidth}{0.4pt}}
\vspace{8pt}

Jedes Mal, wenn du auf Instagram einen Filter anwendest oder in der
Galerie-App ein Foto heller machst, l\"auft im Hintergrund ein
Bildbearbeitungs-Algorithmus. In diesem Projekt programmierst du solche
Filter selbst und siehst sofort, was dein Code mit echten Bildern macht.

% ============================================================================
\section{So funktioniert das Projekt}

\begin{enumerate}
  \item Du arbeitest an deinem eigenen Python-Programm. Du schreibst Filter
    als Funktionen.
  \item Auf einer Website f\"uhrst du ein \textit{Stempelheft}. F\"ur jede
    Aufgabe, die du schaffst, bekommst du einen Stempel.
  \item Mit jedem Stempel wird ein St\"uck eines paint-by-numbers-Bildes auf
    deinem Heft sichtbar. Wer alle Stempel hat, sieht das fertige Bild.
\end{enumerate}

\begin{definitionbox}{Bildfilter}
Ein \term{Bildfilter} ist ein Algorithmus, der ein Bild als Eingabe nimmt
und ein neues Bild zur\"uckgibt. Dabei wird jeder Pixel nach einer
bestimmten Regel ver\"andert.
\end{definitionbox}

% ============================================================================
\section{Spielregeln}

\begin{itemize}
  \item Du arbeitest mit \textbf{verschachtelten Listen}, wie du sie aus
    dem Algorithmen-Skript kennst:
\begin{lstlisting}
bild[zeile][spalte]  # liefert einen Pixel (R, G, B)
\end{lstlisting}
  \item Zum Laden und Speichern von Bilddateien verwendest du die
    Bibliothek \texttt{Pillow}. Setze die folgenden zwei Funktionen
    einmal an den Anfang deiner Datei. Sie laden ein Bild aus dem
    aktuellen Ordner in eine verschachtelte Liste und speichern eine
    verschachtelte Liste wieder als Bild:
\begin{lstlisting}
from PIL import Image

def lade_bild(pfad):
    img = Image.open(pfad).convert("RGB")
    breite, hoehe = img.size
    bild = []
    for y in range(hoehe):
        zeile = []
        for x in range(breite):
            zeile.append(img.getpixel((x, y)))
        bild.append(zeile)
    return bild

def speichere_bild(bild, pfad):
    hoehe = len(bild)
    breite = len(bild[0])
    img = Image.new("RGB", (breite, hoehe))
    for y in range(hoehe):
        for x in range(breite):
            img.putpixel((x, y), bild[y][x])
    img.save(pfad)
\end{lstlisting}
    Aufrufen kannst du sie so:
\begin{lstlisting}
bild = lade_bild("foto.jpg")
speichere_bild(bild, "ergebnis.png")
\end{lstlisting}
  \item Alles andere (jede Filterlogik) schreibst du selbst.
\end{itemize}

\begin{beispiel}{Was ist ein Pixel?}
Ein Pixel wird in Python als Tupel aus drei Zahlen gespeichert, jeweils
zwischen 0 und 255:

\begin{lstlisting}
schwarz = (0, 0, 0)
weiss   = (255, 255, 255)
rot     = (255, 0, 0)
gruen   = (0, 255, 0)
blau    = (0, 0, 255)
gelb    = (255, 255, 0)   # Rot + Gruen
violett = (128, 0, 128)
\end{lstlisting}

Auf einen einzelnen Pixel greifst du so zu:
\begin{lstlisting}
pixel = bild[10][20]   # Pixel in Zeile 10, Spalte 20
r, g, b = pixel        # in seine drei Werte zerlegen
\end{lstlisting}
\end{beispiel}

% ============================================================================
\section{Abgabe}

Du gibst am Ende drei Sachen ab:

\begin{enumerate}
  \item Deine \textbf{Python-Datei(en)} mit allen Filtern.
  \item Drei \textbf{Vorher/Nachher-Bildpaare}: ein Originalfoto und
    daneben jeweils das Ergebnis eines deiner Filter.
  \item Einen \textbf{Bericht} (2 bis 3 A4-Seiten): Was hast du
    gemacht? Welcher Filter hat dich \"uberrascht? Wo bist du steckengeblieben?
\end{enumerate}

Die Abgabe l\"auft \"uber die normale Plattform der Klasse. Wie genau,
besprechen wir gemeinsam.

% ============================================================================
\section{Zeitrahmen}

Das Projekt l\"auft \"uber rund sieben Doppellektionen. Das n\"achste
Kapitel zeigt dir, wie bewertet wird; danach folgt der \textit{Lernpfad}
mit den eigentlichen Aufgaben.

\closechapter

% ============================================================================
% CHAPTER 2: Bewertungskatalog
% ============================================================================
\renewcommand{\chaptertag}{Bildfilter}
\chapter{Bewertungskatalog}

\vspace{-5pt}
{\color{maincolor}\rule{\textwidth}{0.4pt}}
\vspace{4pt}

{\color{maincolor}\textbf{\"Uberblick:}}
\begin{itemize}[leftmargin=*]
  \item Du kannst insgesamt \textbf{120 Punkte} erreichen, verteilt
    auf vier Bereiche.
  \item Den gr\"ossten Teil holst du dir mit den \textbf{Pflichtfiltern}.
    Daneben z\"ahlen dein \textbf{eigener Filter}, deine
    \textbf{Codequalit\"at} und eine ordentliche \textbf{Abgabe}.
  \item Die folgende Grafik zeigt, wieviele Punkte in jedem Bereich zu
    holen sind.
\end{itemize}

\vspace{8pt}

\begin{center}
\begin{tikzpicture}
  \def\W{15}     % Gesamtbreite in cm (entspricht textwidth)
  \def\hgt{1.4}  % Hoehe in cm
  % Anteile aus 120 Punkten: A=50 (0.4167), B=20 (0.1667), C=15 (0.125), D=35 (0.2917)
  \fill[maincolor]    (0,0)              rectangle (\W*0.4167, \hgt);
  \fill[maincolor!72] (\W*0.4167, 0)     rectangle (\W*0.5833, \hgt);
  \fill[maincolor!50] (\W*0.5833, 0)     rectangle (\W*0.7083, \hgt);
  \fill[maincolor!32] (\W*0.7083, 0)     rectangle (\W,        \hgt);
  % Beschriftungen
  \node[white, font=\bfseries]            at (\W*0.2083,\hgt*0.65) {A. Pflichtfilter};
  \node[white, font=\footnotesize]        at (\W*0.2083,\hgt*0.32) {50 Punkte};
  \node[white, font=\bfseries\small]      at (\W*0.5000,\hgt*0.65) {B. Eigen};
  \node[white, font=\scriptsize]          at (\W*0.5000,\hgt*0.32) {20 Pkt.};
  \node[white, font=\bfseries\scriptsize] at (\W*0.6458,\hgt*0.65) {C. Code};
  \node[white, font=\tiny]                at (\W*0.6458,\hgt*0.32) {15 Pkt.};
  \node[white, font=\bfseries]            at (\W*0.8542,\hgt*0.65) {D. Abgabe};
  \node[white, font=\footnotesize]        at (\W*0.8542,\hgt*0.32) {35 Punkte};
\end{tikzpicture}
\end{center}

\vspace{4pt}
{\color{maincolor}\rule{\textwidth}{0.4pt}}
\vspace{8pt}

\renewcommand{\arraystretch}{1.25}

% ============================================================================
\section{A. Pflichtfilter (50 Punkte)}

F\"ur jeden dieser Filter gibt es Punkte, wenn er korrekt funktioniert und
ein sichtbares Ergebnis liefert. Jeder Filter ist eine eigene Funktion.

\begin{center}
\begin{tabularx}{\textwidth}{>{\bfseries}p{4cm} X r}
\toprule
\tableheadercolor
Filter & Was er tut & Pkt. \\
\midrule
Kein Rot
  & Setzt den Rot-Wert aller Pixel auf 0.
  & 5 \\
\midrule
Invertieren
  & Aus jedem Wert wird $255 - $ Wert.
  & 7 \\
\midrule
Graustufen
  & Pro Pixel: \texttt{grau = (R + G + B) // 3}, neuer Pixel
    $(\text{grau}, \text{grau}, \text{grau})$.
  & 7 \\
\midrule
Helligkeit
  & Addiert eine Zahl $h$ auf jeden Farbkanal, klemmt die Werte
    auf $[0, 255]$.
  & 8 \\
\midrule
Schwellwert
  & Pixel \"uber einem Schwellwert $t$ werden weiss, alle anderen
    schwarz.
  & 7 \\
\midrule
Box-Blur
  & Pro Pixel: Durchschnitt aus dem Pixel selbst und seinen 8 Nachbarn
    ($3\times 3$-Fenster).
  & 10 \\
\midrule
Spiegeln
  & Spiegelt das Bild horizontal (linke und rechte Seite tauschen).
  & 6 \\
\bottomrule
\end{tabularx}
\end{center}

% ============================================================================
\section{B. Eigener Filter (20 Punkte)}

Erfinde und implementiere einen \textit{eigenen} Filter, oder w\"ahle
einen aus den Vorschl\"agen am Ende des Lernpfads.

\begin{center}
\begin{tabularx}{\textwidth}{>{\bfseries}p{4cm} X r}
\toprule
\tableheadercolor
Kriterium & Anforderung & Pkt. \\
\midrule
Funktion arbeitet korrekt
  & Der Filter produziert sichtbar das, was er soll.
  & 10 \\
\midrule
Eigenst\"andige Idee
  & Du kannst in zwei S\"atzen erkl\"aren, was dein Filter tut und warum.
  & 5 \\
\midrule
Beispielbild
  & Du zeigst ein Vorher/Nachher-Bild, auf dem der Effekt klar zu sehen ist.
  & 5 \\
\bottomrule
\end{tabularx}
\end{center}

% ============================================================================
\section{C. Codequalit\"at (15 Punkte)}

\begin{center}
\begin{tabularx}{\textwidth}{>{\bfseries}p{4cm} X r}
\toprule
\tableheadercolor
Kriterium & Anforderung & Pkt. \\
\midrule
Aussagekr\"aftige Namen
  & Variablen heissen wie das, was sie sind:
    \texttt{zeile}, \texttt{spalte}, \texttt{neuer\_pixel}, nicht
    \texttt{a}, \texttt{b}, \texttt{tmp}.
  & 5 \\
\midrule
Eine Funktion pro Filter
  & Jeder Filter ist eine eigene Funktion, nicht alles in einem grossen
    Block.
  & 5 \\
\midrule
Clamping korrekt
  & Bei \texttt{helligkeit} bleiben die Werte zwischen 0 und 255.
  & 3 \\
\midrule
Rand bewusst behandelt
  & Beim Box-Blur weisst du, was du am Bildrand machst (z.B.\ einfach
    weglassen).
  & 2 \\
\bottomrule
\end{tabularx}
\end{center}

% ============================================================================
\section{D. Abgabe (35 Punkte)}

\begin{center}
\begin{tabularx}{\textwidth}{>{\bfseries}p{4cm} X r}
\toprule
\tableheadercolor
Kriterium & Anforderung & Pkt. \\
\midrule
Code abgegeben
  & Alle Python-Dateien sind l\"auffahig und an die richtige Stelle hochgeladen.
  & 5 \\
\midrule
Drei Beispielbild-Paare
  & Drei Original/Filter-Paare als Bilddateien.
  & 5 \\
\midrule
Bericht
  & 2 bis 3 A4-Seiten: Was hast du gemacht? Was war schwierig?
    Was hat dich \"uberrascht?
  & 25 \\
\bottomrule
\end{tabularx}
\end{center}

\closechapter

% ============================================================================
% CHAPTER 3: Lernpfad
% ============================================================================
\renewcommand{\chaptertag}{Bildfilter}
\chapter{Lernpfad}

% --- Local color + stamp marker for this chapter and the next ---
\definecolor{cinnabar}{RGB}{177, 56, 46}
\newcommand{\stempel}[2]{%
  \par\vspace{3pt}\noindent
  \begin{tikzpicture}[baseline=-0.3em]
    \node[fill=cinnabar, rounded corners=2pt,
          inner xsep=5pt, inner ysep=2pt,
          font=\scriptsize\bfseries\color{white}] {Stempel #1};
  \end{tikzpicture}\,%
  \textit{\color{cinnabar}#2}\par\vspace{2pt}%
}

\vspace{-5pt}
{\color{maincolor}\rule{\textwidth}{0.4pt}}
\vspace{4pt}

{\color{maincolor}\textbf{Wegweiser:}}
\begin{itemize}[leftmargin=*]
  \item Jeder Schritt baut auf dem vorhergehenden auf.
  \item Jeder Schritt enth\"alt: was du tun sollst, ein kleines Code-Beispiel
    und welchen Stempel du dabei sammelst.
  \item Wenn du eine Station geschafft hast, zeig der Lehrperson dein
    Resultat, und sie setzt den Stempel direkt in dein Heft.
\end{itemize}
\vspace{-2pt}
{\color{maincolor}\rule{\textwidth}{0.4pt}}
\vspace{8pt}

% ============================================================================
\section{Schritt 0: Bild laden}

Bevor du Filter schreiben kannst, musst du ein Bild laden. Dazu
verwenden wir die Python-Bibliothek \texttt{Pillow}. Lege ein Foto in
denselben Ordner wie dein Python-Programm und l\"ade es mit folgendem
Code:

\begin{lstlisting}
from PIL import Image

img = Image.open("foto.jpg").convert("RGB")
breite, hoehe = img.size

bild = []
for y in range(hoehe):
    zeile = []
    for x in range(breite):
        zeile.append(img.getpixel((x, y)))
    bild.append(zeile)

print("Hoehe:", hoehe, " Breite:", breite)
\end{lstlisting}

Die ersten Zeilen \"offnen das Bild und holen seine Gr\"osse. Danach
bauen die zwei verschachtelten Schleifen das Bild als verschachtelte
Liste von Pixel-Tupeln $(R, G, B)$ auf, mit der du im Rest des Projekts
arbeitest. Zum Speichern eines Bildes verwendest du am Ende:

\begin{lstlisting}
ergebnis_img = Image.new("RGB", (breite, hoehe))
for y in range(hoehe):
    for x in range(breite):
        ergebnis_img.putpixel((x, y), bild[y][x])
ergebnis_img.save("ergebnis.png")
\end{lstlisting}

\stempel{1}{Erstes Bild}

% ============================================================================
\section{Schritt 1: Erste Pixel ver\"andern}

Bevor wir richtige Filter schreiben, probieren wir aus, was passiert, wenn
wir einzelne Farbkan\"ale ver\"andern.

\subsection*{Aufw\"armer: Rot weg}

Setze den Rot-Wert aller Pixel auf 0. Was passiert?

\begin{lstlisting}[style=pseudo]
Funktion kein_rot(bild):
    Erstelle ein leeres neues_bild.
    Fuer jede Zeile des Bildes:
        Erstelle eine leere neue_zeile.
        Fuer jeden Pixel (r, g, b) der Zeile:
            Haenge (0, g, b) an neue_zeile an.
        Haenge neue_zeile an neues_bild an.
    Gib neues_bild zurueck.
\end{lstlisting}

\stempel{2}{Kein Rot}

\subsection*{Filter 1: Invertieren}

Beim Invertieren wird jeder Farbwert ``umgedreht'': aus 0 wird 255, aus 255
wird 0, aus 100 wird 155.

\begin{definitionbox}{Invertieren}
F\"ur jeden Pixel $(R, G, B)$:
\[ (R', G', B') = (255 - R,\; 255 - G,\; 255 - B) \]
\end{definitionbox}

\begin{lstlisting}[style=pseudo]
Funktion invertieren(bild):
    Erstelle ein leeres neues_bild.
    Fuer jede Zeile des Bildes:
        Erstelle eine leere neue_zeile.
        Fuer jeden Pixel (r, g, b) der Zeile:
            Haenge (255-r, 255-g, 255-b) an neue_zeile an.
        Haenge neue_zeile an neues_bild an.
    Gib neues_bild zurueck.
\end{lstlisting}

\stempel{3}{Spiegelwelt}

% ============================================================================
\section{Schritt 2: Punktoperationen}

Filter, bei denen jeder neue Pixel nur vom \textit{gleichen} Pixel im
Originalbild abh\"angt, heissen \term{Punktoperationen}.

\subsection*{Filter 2: Graustufen}

In einem Graustufenbild sind alle drei Farbkan\"ale gleich. Der einfachste
Weg: den Durchschnitt der drei Werte nehmen.

\begin{definitionbox}{Graustufen}
F\"ur jeden Pixel $(R, G, B)$:
\[ \text{grau} = \frac{R + G + B}{3}, \qquad
   \text{neuer Pixel} = (\text{grau}, \text{grau}, \text{grau}) \]
\end{definitionbox}

\begin{lstlisting}[style=pseudo]
Funktion graustufen(bild):
    Erstelle ein leeres neues_bild.
    Fuer jede Zeile des Bildes:
        Erstelle eine leere neue_zeile.
        Fuer jeden Pixel (r, g, b) der Zeile:
            Setze grau = (r + g + b) ganzzahlig durch 3.
            Haenge (grau, grau, grau) an neue_zeile an.
        Haenge neue_zeile an neues_bild an.
    Gib neues_bild zurueck.
\end{lstlisting}

\stempel{4}{Graustufen}

\subsection*{Filter 3: Helligkeit}

Um ein Bild heller zu machen, addieren wir auf jeden Farbwert eine Zahl
$h$. Achtung: die Werte m\"ussen zwischen 0 und 255 bleiben! Daf\"ur
schreiben wir eine kleine Hilfsfunktion \texttt{clamp}.

\begin{lstlisting}[style=pseudo]
Funktion clamp(wert):
    Wenn wert kleiner als 0: gib 0 zurueck.
    Wenn wert groesser als 255: gib 255 zurueck.
    Sonst: gib wert zurueck.

Funktion helligkeit(bild, h):
    Erstelle ein leeres neues_bild.
    Fuer jede Zeile des Bildes:
        Erstelle eine leere neue_zeile.
        Fuer jeden Pixel (r, g, b) der Zeile:
            Haenge (clamp(r+h), clamp(g+h), clamp(b+h)) an neue_zeile an.
        Haenge neue_zeile an neues_bild an.
    Gib neues_bild zurueck.
\end{lstlisting}

\stempel{5}{Heller, dunkler}

\begin{aufgabe}{Helligkeit von Hand}
Was kommt heraus, wenn du auf einen Pixel $(200, 100, 50)$ die Helligkeit
\begin{enumerate}
  \item $h = 50$
  \item $h = -150$
  \item $h = 100$ (ohne Clamping)
\end{enumerate}
anwendest?
\end{aufgabe}

% ============================================================================
\section{Schritt 3: Schwellwert}

Beim Schwellwert wird jeder Pixel \textit{entweder} ganz schwarz \textit{oder}
ganz weiss, je nachdem, ob seine Helligkeit \"uber oder unter einem
Schwellwert $t$ liegt.

\begin{definitionbox}{Schwellwert}
F\"ur jeden Pixel berechne zuerst den Grauwert $g$. Dann:
\[ \text{neuer Pixel} = \begin{cases}
   (255, 255, 255) & \text{falls } g \geq t \\
   (0, 0, 0) & \text{sonst}
\end{cases} \]
\end{definitionbox}

\begin{lstlisting}[style=pseudo]
Funktion schwellwert(bild, t):
    Erstelle ein leeres neues_bild.
    Fuer jede Zeile des Bildes:
        Erstelle eine leere neue_zeile.
        Fuer jeden Pixel (r, g, b) der Zeile:
            Setze grau = (r + g + b) ganzzahlig durch 3.
            Wenn grau groesser oder gleich t:
                Haenge (255, 255, 255) an neue_zeile an.
            Sonst:
                Haenge (0, 0, 0) an neue_zeile an.
        Haenge neue_zeile an neues_bild an.
    Gib neues_bild zurueck.
\end{lstlisting}

\stempel{6}{Schwarz-Weiss}

% ============================================================================
\section{Schritt 4: Nachbarpixel und Box-Blur}

Die bisherigen Filter haben jeden Pixel \textit{einzeln} bearbeitet.
Spannender wird es, wenn wir auch die \term{Nachbarpixel} mitnehmen.

\begin{definitionbox}{Box-Blur ($3\times 3$)}
F\"ur jeden Pixel: schau dir den Pixel selbst und seine 8 direkten Nachbarn
an (das $3\times 3$-Fenster). Berechne pro Farbkanal den \textbf{Durchschnitt}
aller 9 Werte. Das ist dein neuer Pixel.

\begin{center}
\begin{tikzpicture}[
  cell/.style={minimum width=10mm, minimum height=10mm, draw=maincolor!50,
    fill=maincolorvlight, font=\scriptsize},
  center/.style={cell, fill=maincolorlight, font=\scriptsize\bfseries},
]
  \node[cell] at (0, 1) {NW};
  \node[cell] at (1.2, 1) {N};
  \node[cell] at (2.4, 1) {NO};
  \node[cell] at (0, 0) {W};
  \node[center] at (1.2, 0) {Pixel};
  \node[cell] at (2.4, 0) {O};
  \node[cell] at (0, -1) {SW};
  \node[cell] at (1.2, -1) {S};
  \node[cell] at (2.4, -1) {SO};
\end{tikzpicture}
\end{center}
\end{definitionbox}

\begin{lstlisting}[style=pseudo]
Funktion box_blur(bild):
    Erstelle ein leeres neues_bild.
    Fuer jede Zeile von 1 bis hoehe-2:
        Erstelle eine leere neue_zeile.
        Fuer jede Spalte von 1 bis breite-2:
            Setze summe_r, summe_g, summe_b auf 0.
            Fuer dz in {-1, 0, 1}:
                Fuer ds in {-1, 0, 1}:
                    Lies Pixel (r, g, b) an (zeile+dz, spalte+ds).
                    summe_r = summe_r + r
                    summe_g = summe_g + g
                    summe_b = summe_b + b
            Haenge (summe_r//9, summe_g//9, summe_b//9) an neue_zeile an.
        Haenge neue_zeile an neues_bild an.
    Gib neues_bild zurueck.
\end{lstlisting}

Beachte: wir beginnen bei Zeile/Spalte~1 und h\"oren eine Zeile/Spalte vor
dem Ende auf. So haben wir immer ein vollst\"andiges $3\times 3$-Fenster.
Die \"aussersten Pixel fallen dabei weg. Das ist eine bewusste Entscheidung
und f\"ur unsere Zwecke v\"ollig in Ordnung.

\stempel{7}{Nachbarn mischen}

\begin{aufgabe}{Mehrfach blurren}
Was passiert, wenn du den Box-Blur \textit{mehrfach hintereinander} auf
dasselbe Bild anwendest? Probiere 1$\times$, 5$\times$ und 20$\times$ aus.
\end{aufgabe}

% ============================================================================
\section{Schritt 5: Bild umordnen, Spiegeln}

Bisher haben wir Pixel \textit{ver\"andert}. Jetzt ordnen wir sie um. Beim
horizontalen Spiegeln tauschen wir die linke und die rechte Seite.

\begin{lstlisting}[style=pseudo]
Funktion spiegeln(bild):
    Erstelle ein leeres neues_bild.
    Fuer jede Zeile des Bildes:
        Erstelle eine leere neue_zeile.
        Fuer jede Spalte von 0 bis breite-1:
            Haenge bild[zeile][breite-1-spalte] an neue_zeile an.
        Haenge neue_zeile an neues_bild an.
    Gib neues_bild zurueck.
\end{lstlisting}

\stempel{8}{Spiegelei}

% ============================================================================
\section{Schritt 6: Dein eigener Filter}

Jetzt bist du dran! Erfinde \textit{einen} eigenen Filter. Hier ein paar
Ideen (du darfst auch etwas ganz Eigenes machen):

\begin{itemize}
  \item \textbf{Vertikal spiegeln:} obere und untere H\"alfte tauschen.
  \item \textbf{Posterisierung:} runde jeden Kanal auf Vielfache von 64
    (Werte werden dann nur noch 0, 64, 128, 192).
  \item \textbf{Verpixelung:} ersetze jeden $k \times k$-Block durch
    die Durchschnittsfarbe dieses Blocks.
  \item \textbf{Einfache Kantenerkennung:} f\"ur jeden Pixel den Unterschied
    zum rechten Nachbarn berechnen.
  \item \textbf{Nur ein Kanal:} alle Farben ausser Rot, Gr\"un oder Blau
    auf 0 setzen.
\end{itemize}

\stempel{9}{Eigener Filter}

% ============================================================================
\section{Schritt 7: Abgabe}

Bereite die Abgabe vor:

\begin{enumerate}
  \item Pr\"ufe, dass alle Pflichtfilter laufen.
  \item Suche drei Originalbilder aus und wende je einen Filter darauf an.
    Das ergibt drei \textit{Vorher/Nachher}-Paare.
  \item Schreibe einen Bericht von 2 bis 3 A4-Seiten. Was
    war einfach? Was war schwierig? Welcher Filter hat dich \"uberrascht?
  \item Gib Code, Bilder und Bericht ab.
\end{enumerate}

\stempel{10}{Auf dem Gipfel}

\closechapter

% ============================================================================
% CHAPTER 4: Stempel-Wanderung
% ============================================================================
\renewcommand{\chaptertag}{Bildfilter}
\chapter{Stempel-Wanderung}

% cinnabar color and \stempel command are defined in lernpfad.tex (loaded earlier).

\vspace{-5pt}
{\color{maincolor}\rule{\textwidth}{0.4pt}}
\vspace{4pt}

{\color{maincolor}\textbf{Worum es geht:}}
\begin{itemize}[leftmargin=*]
  \item Auf japanischen Wanderwegen gibt es an jeder Station einen Stempel,
    den du in dein Heft dr\"uckst. Wer alle Stempel hat, hat den Berg geschafft.
  \item Wir machen das mit Code: jede gemeisterte Station gibt
    \textbf{einen Stempel}. Mit jedem Stempel wird ein Teil deines
    \textit{paint-by-numbers}-Bildes sichtbar.
  \item Es gibt insgesamt \textbf{10 Stationen}, vom ersten geladenen
    Pixel bis zur fertigen Abgabe.
\end{itemize}
\vspace{-2pt}
{\color{maincolor}\rule{\textwidth}{0.4pt}}
\vspace{8pt}

% ============================================================================
\section{So funktioniert es}

\begin{enumerate}
  \item Zu Beginn des Projekts bekommst du eine \textbf{Login-Karte} mit
    deinem Benutzernamen und einem Passwort.
  \item Damit \"offnest du die Website \textit{Bildfilter Stempel-Wanderung}
    im Browser und meldest dich an.
  \item Wenn du eine Station geschafft hast, zeigst du dein Resultat der
    Lehrperson.
  \item Die Lehrperson setzt dir den Stempel direkt in dein Heft. Die
    n\"achsten zwei Farben auf deinem Bild werden sofort sichtbar.
\end{enumerate}

\begin{definitionbox}{Login-Karte}
Auf deiner Login-Karte stehen:
\begin{itemize}
  \item dein \textbf{Benutzername} (in der Regel dein Vorname, klein geschrieben),
  \item dein \textbf{Passwort} (sechs Zeichen),
  \item die \textbf{Adresse} der Website.
\end{itemize}
Bewahre die Karte gut auf. Wenn du sie verlierst, sag der Lehrperson Bescheid,
dann gibt es ein neues Passwort.
\end{definitionbox}

% ============================================================================
\section{Der Trail}

\begin{center}
\begin{tikzpicture}[
  station/.style={
    circle, draw=cinnabar, fill=cinnabar!8,
    minimum size=10mm, font=\small\bfseries\color{cinnabar},
    line width=0.9pt, inner sep=0pt,
  },
  trail/.style={cinnabar!55, line width=1.4pt, rounded corners=4pt},
  caption/.style={font=\scriptsize, align=center, color=black!75, text width=24mm},
]
  % Row 1 (left to right)
  \node[station] (s1) at (0, 0)  {1};
  \node[station] (s2) at (3, 0)  {2};
  \node[station] (s3) at (6, 0)  {3};
  \node[station] (s4) at (9, 0)  {4};
  % Row 2 (right to left)
  \node[station] (s5) at (9, -2.5) {5};
  \node[station] (s6) at (6, -2.5) {6};
  \node[station] (s7) at (3, -2.5) {7};
  % Row 3 (left to right)
  \node[station] (s8) at (3, -5)   {8};
  \node[station] (s9) at (6, -5)   {9};
  \node[station] (s10) at (9, -5)  {10};

  % Trail (winding)
  \draw[trail] (s1) -- (s2) -- (s3) -- (s4)
               -- (s5) -- (s6) -- (s7)
               -- (s8) -- (s9) -- (s10);

  % Captions
  \node[caption] at (0, 0.85)   {Erstes\\Bild};
  \node[caption] at (3, 0.85)   {Kein Rot};
  \node[caption] at (6, 0.85)   {Spiegelwelt};
  \node[caption] at (9, 0.85)   {Graustufen};
  \node[caption] at (9, -1.65)  {Heller,\\dunkler};
  \node[caption] at (6, -1.65)  {Schwarz-\\Weiss};
  \node[caption] at (3, -1.65)  {Nachbarn\\mischen};
  \node[caption] at (3, -4.15)  {Spiegelei};
  \node[caption] at (6, -4.15)  {Eigener\\Filter};
  \node[caption] at (9, -4.15)  {Auf dem\\Gipfel};

  % Decorative summit marker after last station
  \node[font=\Large\color{cinnabar}] at (10.4, -5) {$\blacktriangle$};
\end{tikzpicture}
\end{center}

\vspace{6pt}

% ============================================================================
\section{Stationen im Detail}

\renewcommand{\arraystretch}{1.2}

\begin{center}
\begin{tabularx}{\textwidth}{>{\bfseries\color{cinnabar}}c >{\bfseries}p{3.6cm} X p{1.8cm}}
\toprule
\tableheadercolor
\# & Stempel & Anforderung & Lernpfad \\
\midrule
1  & Erstes Bild
   & Bild mit \texttt{lade\_bild} in eine Liste laden, H\"ohe und Breite ausgeben.
   & Schritt 0 \\
2  & Kein Rot
   & Rot-Wert aller Pixel auf $0$ setzen und speichern.
   & Schritt 1 \\
3  & Spiegelwelt
   & \texttt{invertieren(bild)}: jeder Wert wird zu $255 -$ Wert.
   & Schritt 1 \\
4  & Graustufen
   & \texttt{graustufen(bild)}: Mittelwert $\frac{R+G+B}{3}$.
   & Schritt 2 \\
5  & Heller, dunkler
   & \texttt{helligkeit(bild, h)} mit Clamping auf $[0, 255]$.
   & Schritt 2 \\
6  & Schwarz-Weiss
   & \texttt{schwellwert(bild, t)}: Pixel \"uber $t$ werden weiss.
   & Schritt 3 \\
7  & Nachbarn mischen
   & \texttt{box\_blur(bild)}: $3\times 3$-Nachbarschaft mitteln.
   & Schritt 4 \\
8  & Spiegelei
   & \texttt{spiegeln(bild)}: horizontal spiegeln.
   & Schritt 5 \\
9  & Eigener Filter
   & Ein selbst erfundener Filter, der sichtbar funktioniert.
   & Schritt 6 \\
10 & Auf dem Gipfel
   & Code, drei Vorher/Nachher-Bildpaare und Bericht (2 bis 3 A4-Seiten) abgegeben.
   & Schritt 7 \\
\bottomrule
\end{tabularx}
\end{center}

% ============================================================================
\section{Das Bild}

Dein Stempelheft zeigt eine \textit{paint-by-numbers}-Vorlage des Berges
Fuji mit einer japanischen Pagode. Am Anfang ist nur die Vorlage mit den
Zahlen zu sehen. Mit jedem Stempel wird ein St\"uck der Vorlage durch das
\textit{fertige Bild} ersetzt, das, was du am Ende der Wanderung in
voller Farbe siehst.

Insgesamt ist das Bild in zehn Bereiche aufgeteilt, einer pro Stempel.

\closechapter


\end{document}
